\chapter{IMPLEMENTATION}
\label{ch:implementation}

\section{Introduction}

This chapter describes the way in which the design of Chapter~\ref{ch:methodology}
is realised in software, that is how the LabelMe annotations are converted into a
trainable format, how the instrument banner and the physical calibration are
handled, how augmentation is applied after partitioning, and how the six methods are
configured and trained.

\section{System Architecture}

The pipeline is organised as five stages, each independently runnable so the
expensive stage, training, does not have to be repeated when only the reporting
changes. Ingest reads the LabelMe JSON files, recovers specimen and
magnification and extracts the pixel calibration from the instrument banner.
Prepare crops the banner, rasterises the polygons into instance masks and emits a
COCO-format annotation file. Partition assigns images to the four
leave-one-specimen-out folds. Train and predict fits each method on each fold's
training partition and predicts on its held-out specimen. Evaluate and report
computes segmentation metrics, counting and merge/split statistics and converts
detections into physical units.

\section{Data Preparation}

\subsection{Annotation Parsing and Rasterisation}
\label{sec:parsing}

Each LabelMe file stores a list of shapes, each carrying a class label, a shape
type and a list of vertices. All 606 shapes are polygons carrying the label
\texttt{bead}, rasterised into per-instance binary masks at full resolution.
Object area is computed analytically from the polygon vertices using the shoelace
formula rather than by counting mask pixels, which avoids a rasterisation bias
significant for the smallest beads, since at $500\times$ the median bead is only
about 14 pixels across, so a one-pixel boundary error changes its area by roughly
15\%. The equivalent circular diameter follows as $d = 2\sqrt{A/\pi}$ and is
converted to micrometres using the magnification-specific pixel size from
Table~\ref{tab:calibration}.

The polygons are kept as polygons rather than flattened into a single label map.
Measurement shows that 8.8\% of all annotated bead pixels are claimed by more than
one polygon, and a label map assigns each pixel to exactly one object, so it would
discard that overlap silently, and precisely in the regions where merge and split
errors are decided. Instance masks are therefore rasterised on demand, one plane
per object.

\subsection{Banner Handling and Calibration}

The instrument banner occupies the bottom 32 rows of every image. Calibration
is performed before cropping, where the scale bar is located as the longest run of
saturated pixels in the lower-left region of the banner and its pixel length divided
by the physical length printed beside it gives the pixel size. The recovered values
are cross-checked by verifying that the field width scales inversely with
magnification, as reported in Section~\ref{sec:calibration}.

The banner rows are then removed, giving a $1024 \times 736$ working image. The
polygon coordinates require no adjustment, since the crop is taken from the bottom
edge and no annotation extends into the banner region.

\subsection{Fold Assignment}

Fold membership is derived from the specimen prefix of the filename, so all
images of one specimen move together. The four folds are S1, S2, S3 and S4. The
assignment is written to disk as an explicit manifest rather than recomputed at
training time, so every method is guaranteed to see identical partitions and
the comparison between them is exact.

The deliberately flawed control configuration described in
Section~\ref{sec:protocol} uses the same machinery with the specimen constraint
disabled, splitting the 11 images at random.

\section{Augmentation}

Augmentation is applied on the fly during training, after partitioning, for the
reasons set out in Section~\ref{sec:supplied_aug}. The pre-computed copies shipped
with the dataset are not used.

The transformations are chosen so as to be label-preserving for this modality.
Images of a non-woven membrane have no canonical orientation, so horizontal and
vertical flips and rotations in multiples of $90^{\circ}$ are physically meaningful.
Random scaling is included deliberately and with a wider range than would be usual,
since scale robustness across the $500\times$ to $3000\times$ span is the central
difficulty of the dataset. Brightness and contrast jitter model the variation in
detector gain between sessions. Elastic and shear deformations are excluded, as they
would distort bead shape, and bead shape is a measured output rather than an
incidental property.

Random cropping to fixed-size tiles is used both as augmentation and as a practical
measure, since a $1024 \times 736$ image containing up to 152 small beads is
awkward to fit in memory at the batch sizes needed for stable training.

\section{Model Configuration}

\subsection{Mask R-CNN}

The primary model uses a ResNet-50 backbone with a feature pyramid network,
initialised from COCO-pretrained weights. Two departures from the defaults are
required by the data.

The anchor sizes are set from the empirical bead size distribution rather than left
at their COCO values. The median bead occupies 0.04\% of the image area, which is
far smaller than a typical COCO object, and default anchors would leave the smallest
beads without a well-matched proposal.

The maximum number of detections per image is raised well above the COCO default of
100, since a single $500\times$ image contains up to 152 beads and the default
would silently truncate the count, which is an error that would corrupt precisely
the quantity being measured.

\subsection{U-Net Baseline}

The semantic baseline uses a ResNet encoder initialised from ImageNet weights
together with a symmetric decoder with skip connections. It is trained with a
combined cross-entropy and Dice loss to counter the strong
foreground/background imbalance, since bead pixels constitute only 2 to 20\% of a
image. Writing $\hat{y}_i$ for the predicted probability at pixel $i$ and
$y_i \in \{0,1\}$ for the annotation,
%
\begin{equation}
    \mathcal{L} =
    \underbrace{-\frac{1}{N}\sum_{i=1}^{N}
      \bigl[\, y_i \log \hat{y}_i + (1 - y_i) \log (1 - \hat{y}_i) \,\bigr]}
      _{\textstyle \text{cross-entropy}}
    \;+\;
    \underbrace{1 - \frac{2 \sum_i \hat{y}_i y_i + \varepsilon}
                        {\sum_i \hat{y}_i + \sum_i y_i + \varepsilon}}
      _{\textstyle \text{Dice}},
    \label{eq:dice_bce}
\end{equation}
%
with $\varepsilon = 1$ guarding the degenerate case of an empty crop.
Cross-entropy alone is minimised well by a nearly empty prediction when the
foreground is this sparse, whereas the Dice term is scale-invariant in the number of
foreground pixels and supplies the gradient that prevents this.

The instances are recovered from the predicted mask by connected-component
labelling, with no boundary-aware post-processing. The omission is deliberate, since
connected components cannot separate two beads whose predicted regions touch, and
adding a boundary-aware correction would hide that limitation rather than measure
it. Whether it costs anything on this data is reported as the merge rate in
Chapter~\ref{ch:results}.

\subsection{Faster R-CNN}

Faster R-CNN is built from the same torchvision family as Mask R-CNN and from the
same COCO-pretrained ResNet-50 and feature pyramid weights, with the mask predictor
absent. The anchor ladder, the cap of 400 detections per image, the region proposal limits,
the optimiser, the learning-rate schedule and the input-size pinning are set by the
same code path, so the two models differ in the mask branch and in nothing else
that was configured.

The consequence for the output format is deliberate. The predictions are written as
boxes with no segmentation field, and the evaluation module therefore leaves every
mask-level metric unset for this method rather than filling each box with a
rectangle to score it. ``Did not segment'' and ``segmented badly'' are
different claims, and a zero would report the second while meaning the first.

\subsection{YOLOv8 and YOLOv5u}

YOLOv8 is trained through the Ultralytics package rather than through the
torchvision path, which imposes two implementation differences on which the
results depend.

The package expects a dataset on disk in its own layout, so each fold is exported
to a scratch directory as images plus one label file per image, with the instance
polygons converted to normalised vertex lists. That export is written to local
disk rather than to the mounted Drive, since the iteration count converts to roughly
750 epochs over eight images and the package rewrites its \texttt{last.pt}
checkpoint at the end of every epoch, which over a network filesystem would
dominate the run time.

The package also drives its own training loop, so the schedule, the augmentation
and the optimiser are its defaults rather than those in
Table~\ref{tab:training_config}, with three exceptions set explicitly: the input
size at 1024, mosaic augmentation disabled and the detection cap at 400. Its
predictions are converted to the same run-length-encoded mask format every other
mask-predicting method writes.

Both YOLO families are handled by one script, which selects the family from the
weights filename and writes each to its own prediction file. Both are reported in
Chapter~\ref{ch:results}, although not on the same footing, since the package
ships no segmentation variant of YOLOv5 and that family therefore enters on the
box-only path alongside Faster R-CNN.

One point of nomenclature has to be stated, since the version numbers invite a
reading the weights do not support. The checkpoint used is \texttt{yolov5su.pt},
which is not the YOLOv5 of the original release: Ultralytics distributes the
YOLOv5 backbone fitted with the anchor-free, objectness-free split head of YOLOv8.
Comparing it against \texttt{yolov8s-seg} therefore varies mostly the backbone and
the presence of a mask branch, and is not the anchor-based against anchor-free
comparison the version numbers suggest. The results are labelled YOLOv5u
throughout.

\subsection{Classical Baseline}

The classical pipeline is implemented with scikit-image, comprising CLAHE contrast
normalisation, Otsu thresholding, morphological opening and closing, the distance
transform, local-maximum markers and watershed. Its three free parameters, namely
the morphological kernel size, the minimum accepted object area in $\mu$m$^2$ and
the marker suppression distance, are tuned by grid search on the training partition
of each fold and are then held fixed for that fold's test specimen, so it is
subject to the same protocol as the learned methods.

\section{Training Setup}

Each of the four folds is trained independently from the same initialisation, and no
fold's held-out specimen influences any choice made during its training, neither the
number of iterations nor the score threshold applied to the output. The threshold is
selected separately for each fold by maximising $F_1$ on that fold's own training
partition.

The configuration is reported in Chapter~\ref{ch:results},
Table~\ref{tab:training_config}, rather than here, because the values that matter
are the ones that the runs actually used. Stating a planned configuration in this
chapter and a measured one in the next invites the two to drift apart, and it is the
measured one against which the results have to be read.

\section{Development Environment}
\label{sec:environment}

Development is in Python with PyTorch \cite{paszke2019pytorch} and torchvision.
Annotation handling and figures use NumPy, Pillow and Matplotlib, the classical
baseline uses scikit-image \cite{vanderwalt2014scikit} and SciPy, and average
precision is computed with the reference COCO implementation
\cite{lin2014microsoft}.

No local GPU was available, so the training stages run on Google Colab with an
NVIDIA T4. This constrains the design in two ways. A session can be interrupted,
so each stage writes its predictions to persistent storage and can be resumed
independently rather than the pipeline being one long run. And the total GPU time
is limited, which is why the comparison is between six methods under one protocol
rather than a wider hyperparameter search, since the leave-one-specimen-out
protocol already multiplies every configuration by four.

The iteration counts and every other setting the runs actually used are recorded
in Table~\ref{tab:training_config}, which \texttt{make\_tables.py} writes from the
metadata each training script stores alongside its predictions. Wall-clock cost is
instrumented separately by \texttt{runtime.py} and reported in
Table~\ref{tab:runtime}.

\section{Challenges and Solutions}
\label{sec:challenges}

Six problems were encountered while building the pipeline, and each changed a
design decision.

The augmented copies could not be used as supplied. The dataset arrives
as 143 image--annotation pairs, which initially appears to be a usable sample
size, while the filenames reveal 11 originals with 12 transformed copies each.
The transformations are correctly applied, so the copies are not defective, but
they are not independent either, and letting them cross a partition boundary would
have produced optimistic results.

The banner is augmented along with the image. The flipped copies show the
instrument banner mirrored and the translated copies carry black padding along the
vacated edge, which reinforced the decision to crop it before any processing.

Object size varies sixfold across the dataset. The median bead is 14
pixels across at $500\times$ and 76 at $3000\times$, which drove the choice of a
feature pyramid backbone, the resetting of the anchor scales and the wide
random-scaling range in the augmentation policy.

Augmentation initially shifted every mask by one pixel. Reflections were
first implemented as $x \mapsto W - 1 - x$, correct for an integer pixel index but
not for polygon vertices, which are continuous coordinates and reflect as
$x \mapsto W - x$. Half a pixel at each edge is invisible on a large object, but
the mean bead here is 17 pixels across, so a one-pixel boundary shift moves
roughly 13\% of its area on every flipped training sample. The fault was found by
rasterising the transformed polygons and comparing them against the directly
transformed raster mask, which agree to an intersection over union of 0.9995 under
the correct convention and 0.87 under the incorrect one. That check is now
permanent in the test suite, since the failure produced no error message and no
visibly wrong image, only slightly wrong labels.

Evaluating a thousand objects per image was forty times slower than
necessary. Instances were initially held as full-frame masks, one
$1024 \times 736$ array per object, and the classical baseline emits over a
thousand objects for a single image, so its parameter search spent almost all
its time reading zeros: a single Aggregated Jaccard Index computation took 5.95
seconds. Holding each instance as a bounding box plus a mask covering only that
box, and computing intersections only for pairs whose boxes overlap, reduced that
to 0.02 seconds and the full parameter search from over an hour to a few minutes.

The smoke run destroyed a completed set of results. The fast correctness
check of Section~\ref{sec:deployment} originally wrote its predictions to the same
directory as a real run, so re-running it after a completed four-fold experiment
silently overwrote several hours of predictions with output from eight training
iterations. A smoke run produces a well-formed predictions file, so nothing fails
and nothing warns, and the only symptom is that every subsequently reported number
is wrong. The fix is structural rather than procedural: a smoke run now redirects
its working and results directories before any pipeline module is imported.

\section{Code Organization}

Table~\ref{tab:modules} lists the modules and what each is responsible for.

\begin{table}[htbp]
    \centering
    \caption{The modules of the implementation. Each concern is owned by one
    file, so a rule such as the partitioning guarantee is written once and
    inherited everywhere it applies.}
    \label{tab:modules}
    \begin{tabular}{p{0.32\textwidth}p{0.58\textwidth}}
        \toprule
        Module & Responsibility \\
        \midrule
        \texttt{config.py} & Dataset facts, calibration constants and paths, in one place so no constant is defined twice. \\
        \texttt{dataset\_stats.py} & Descriptive statistics and the dataset figures of Chapter~\ref{ch:methodology}. \\
        \texttt{prepare\_data.py} & Banner cropping and COCO conversion. \\
        \texttt{folds.py} & The fold manifests for both protocols. \\
        \texttt{data\_io.py} & Annotation loading, polygon rasterisation and the augmentation geometry, free of any deep-learning dependency so the classical baseline and the evaluation code run without PyTorch. \\
        \texttt{datasets.py}, \texttt{models.py} & The PyTorch dataset, the Mask R-CNN configuration and the U-Net definition. \\
        \texttt{train\_maskrcnn.py}, \texttt{train\_fasterrcnn.py}, \texttt{train\_unet.py}, \texttt{train\_yolo.py}, \texttt{classical.py} & The six methods, the two one-stage models being driven by a single script that selects the family from the weights filename. \\
        \texttt{predio.py} & The shared prediction format. Every method writes COCO results, the mask-predicting ones as run-length-encoded masks and the box-only ones as boxes with no segmentation field. \\
        \texttt{metrics.py}, \texttt{evaluate.py} & Segmentation metrics, counting and merge/split statistics, and physical quantification. \\
        \texttt{runtime.py} & The wall-clock instrumentation behind Table~\ref{tab:runtime}. \\
        \texttt{make\_tables.py}, \texttt{make\_figures.py} & The tables and figures of Chapter~\ref{ch:results} and the dataset figures, generated from the measured metrics. \\
        \texttt{run\_all.py} & Runs the pipeline end to end, with a \texttt{-{}-smoke} mode that executes a few iterations to verify nothing crashes before committing hours of GPU time. \\
        \texttt{notebooks/Thesis\_Colab.ipynb} & The notebook used to run the training stages. \\
        \bottomrule
    \end{tabular}
\end{table}

Chapter~\ref{ch:results} is populated by \texttt{make\_tables.py} rather than by
hand, so no number can reach the thesis without a run having produced it.

\section{Deployment}
\label{sec:deployment}

The pipeline is deployed as a Google Colab notebook that clones the repository,
mounts the specimen data from Drive and runs each stage as a separate cell. Three
properties of that arrangement are deliberate.

Each stage is independently resumable. A Colab session can be
disconnected at any point, and a four-fold run of six methods takes long enough
that this is likely rather than merely possible, so every stage writes its
predictions to persistent storage before the next begins.

A smoke path runs the entire pipeline in about two minutes. It executes
eight training iterations per fold and a single classical parameter combination.
Its numbers are meaningless; its purpose is to establish that nothing crashes
before committing hours of GPU time. It writes to a separate working directory,
for the reason given in Section~\ref{sec:challenges}.

Reporting is reproducible from stored metrics. The evaluation and
table-generation stages read only the saved predictions, so the tables and figures
of Chapter~\ref{ch:results} regenerate without repeating any training, and a
change in how a result is presented cannot alter the result.

The complete repository, including the notebook, is available at
\url{https://github.com/ErlindSkura/Erlind_Skura_Thesis}.

The pipeline separates ingest, preparation, partitioning, training and evaluation
so the partitioning guarantee of Chapter~\ref{ch:methodology} is enforced in one
place and inherited by every method. The two decisions that most affect the
validity of the results, discarding the pre-computed augmentations and cropping
the instrument banner, both follow from properties of the data discovered during
preparation rather than from convention. The next chapter reports what the
pipeline produces.
